\worksheettitle
  {Домашняя практика: ответы}
  {\modulemark{M02-H} Измеряем и строим отрезки}

\begin{teacherbox}[Критерии]
  В задачах со сдвинутым началом обязательна разность показаний. В построении
  проверяются два отмеченных конца и длина при измерении линейкой.
\end{teacherbox}

\section{Ответы}

\begin{enumerate}
  \item a) $6\text{ см }2\text{ мм}$; б)
    $7\text{ см }5\text{ мм}-1\text{ см }7\text{ мм}
      =5\text{ см }8\text{ мм}$; в)
    $9\text{ см}-3\text{ см }4\text{ мм}=5\text{ см }6\text{ мм}$.
  \item $AB=4\text{ см }3\text{ мм}$, $CD=6\text{ см }1\text{ мм}$.
  \item Да: $5\text{ см }2\text{ мм}=52\text{ мм}$.
  \item Причина: показание конца принято за длину. Верно:
    $8\text{ см }4\text{ мм}-2\text{ см}=6\text{ см }4\text{ мм}$.
\end{enumerate}

\begin{resultbox}[Маршрут]
  Повторяющаяся подмена длины показанием требует M02-R. Единичная ошибка в
  переводе сантиметров и миллиметров может потребовать дополнительной опоры
  перед M03.
\end{resultbox}
